% !PRINTED 103
\weberpage{103}

\noindent Eine Function $n$ten Grades lässt sich in $n$ lineare Factoren zerlegen.
\medskip

Die Grössen $\alpha,\beta\ldots\nu$ sind dann alle Wurzeln von $f(x)$ und ihre Zahl ist also $n$. Es ist aber nicht ausgeschlossen, dass unter den $\alpha,\beta\ldots\nu$ dieselbe Zahl mehrmals vorkommt. Dann würde $f(x)$ weniger als $n$ Wurzeln haben, während doch der Satz bestehen bleibt, dass $f(x)$ in $n$ lineare Factoren zerlegbar ist. Um die Uebereinstimmung herzustellen, ist man übereingekommen, wenn $x-\alpha$ mehrmals in $f(x)$ aufgeht, $\alpha$ unter den Wurzeln mehrfach zu zählen und also von einfachen, zweifachen, dreifachen etc. Wurzeln zu sprechen.

Nach dem Begriff der derivirten Functionen [§. 13 (2)] ist, wenn $\alpha$ eine beliebige Grösse ist,
\begin{equation}
 f(x)=f(\alpha)+(x-\alpha)f'(\alpha)+\frac{(x-\alpha)^2}{1\cdot2}f''(\alpha)
 +\frac{(x-\alpha)^3}{1\cdot2\cdot3}f'''(\alpha)+\cdots .
\tag{4}
\end{equation}
Wird also nun wieder angenommen, dass $f(\alpha)$ verschwindet, so ergiebt sich
\[
 f_1(x)=f'(\alpha)+\frac{x-\alpha}{1\cdot2}f''(\alpha)+\frac{(x-\alpha)^2}{1\cdot2\cdot3}f'''(\alpha)+\cdots,
\]
also
\[
 f_1(\alpha)=f'(\alpha).
\]
Es ist also $\alpha$ eine Doppelwurzel von $f(x)$, wenn mit $f(\alpha)$ gleichzeitig $f'(\alpha)$ verschwindet. Auch die Formel (4) zeigt, dass nur unter dieser Voraussetzung $f(x)$ durch $(x-\alpha)^2$ theilbar ist. Diese Schlussweise lässt sich weiter ausdehnen und führt zu dem Satze:

\medskip
\noindent III. \gesperrt{Die nothwendige und hinreichende Bedingung dafür, dass $\alpha$ eine $m$-fache Wurzel von $f(x)$ ist, ist die, dass $f(\alpha),f'(\alpha),f''(\alpha),\ldots f^{(m-1)}(\alpha)$ verschwinden, $f^{(m)}(\alpha)$ nicht verschwindet.}

\begin{center}
{\large \S.\ 31.}

\vspace{0.5em}
{\large\gesperrt{Stetigkeit ganzer Functionen.}}
\end{center}

Wenn, wie bisher,
\[
 f(x)=a_0x^n+a_1x^{n-1}+\cdots+a_n
\]
eine ganze Function von $x$ ist, so beweisen wir zunächst folgenden Satz:

\medskip
